% MT001 / explanation — canonical editable content.
逻辑回归先算实数分数\(f(x)=w^Tx+b\)，再用sigmoid函数\(\sigma(f)=1/(1+e^{-f})\)把它变成0到1之间的正类概率；另一类的概率是\(1-\sigma(f)\)。训练通过调整\(w,b\)，使正确标签的概率更高；学习率决定每次调整迈多大一步，正则强度决定我们多不愿意使用大权重。


在附近小幅改变参数，损失都不会更低，称为局部最小；在所有可能的参数中损失最低，称为全局最小。凸函数的图像不高于任意两点间的连线，因此它的局部最小值也是全局最小值。不过，最低位置可以不止一个。


补充例子：\(L(w_1,w_2)=w_1^2\)不依赖\(w_2\)。只要\(w_1=0\)，无论\(w_2\)取0、2还是其他值，损失都是最小值0。这是凸目标，却有一整条同样低的“平底”。



\begin{itemize}

\item \textbf{A按原解视为正确，但表述有条件。} 标准线性LR的负log似然是凸函数，因而没有损失更高的局部最小值；这不保证最优参数唯一，也不保证存在有限的最优权重。原解意在说明前一个性质。

\item \textbf{B错误，选。} 梯度下降写成\(w_{t+1}=w_t-\eta_t\nabla L(w_t)\)。这里\(t\)是第几次更新，\(L\)是训练损失；梯度\(\nabla L\)把各个权重的局部变化率列在一起，指出损失增加最快的方向。沿反方向走足够小的一步，损失会下降；步长太大则可能越过低点。\(\eta_t\)控制这一步的大小，可以固定，也可以随训练调整，没有“必须固定”的要求。

\item \textbf{C错误，选。} 常见目标为\(\|w\|^2/C+\sum_i\ell_i\)，增大\(C\)使权重惩罚变弱。模型更能贴合训练数据，但测试表现可能改善，也可能过拟合后变差；连分类错误率也不必随参数连续、严格变化。

\item \textbf{D正确。} 训练使用连续概率，便于计算损失和梯度；测试只需要类别，可以判断\(f(x)>0\)。因为\(\sigma(0)=0.5\)且sigmoid递增，分数取正与概率超过0.5等价。

\end{itemize}


选读：无正则且数据线性可分时，把已正确分类的权重越放越大，真实标签的概率就越接近1，损失趋向0，却可能始终不在任何有限权重处等于0。这说明“目标凸”和“必有唯一有限解”是两回事。考试按配套答案选B、C，并保留A的精确条件。


\textbf{出处：} 2020B期中Quiz Q1，原卷第2页（5分）


\textbf{答案依据：} 2020B期中 Q1 随卷／配套参考答案；条件与解释经复核。
